% TEMPLATE-MILC-v3.tex - plantilla maestra UNICA de las guias MILC v3.
%
% La usan las tres familias de guias, cada una con su driver:
%   · Grados 6-11  -> scripts/build-guias-g11.py
%   · Bebras       -> scripts/build-guias-bebras.py
%   · Semillero    -> scripts/build-guias-semillero.py
%
% Antes habia tres copias casi identicas (~400 lineas iguales, 6 distintas):
% cada mejora habia que aplicarla tres veces, y en la practica no se hacia
% (el motor de recursos vivia solo en dos de ellas). Lo que varia entre
% familias esta parametrizado en 6 placeholders que cada driver llena
% (se nombran aqui SIN su sintaxis real: el builder reemplaza por texto plano,
% tambien dentro de los comentarios, y un valor multilinea rompe el preambulo):
%
%   HEADER_IZQ        encabezado izquierdo de pagina
%   HEADER_DER        encabezado derecho de pagina
%   PORTADA_ETIQUETA  rotulo sobre el numero grande (GRADO/MOMENTO/MODULO)
%   PORTADA_SERIE     linea de serie de la portada
%   PORTADA_ID        identificador de la guia en la portada
%   PORTADA_CIRCULO   el circulito de grado (°), o vacio si no aplica
%
% El resto de claves las documenta content/guias/_SCHEMA.md. El script valida
% que no queden placeholders sin reemplazar antes de escribir el archivo final.

\documentclass[11pt,letterpaper]{article}
\usepackage[margin=1.75cm]{geometry}
\usepackage[table]{xcolor}
\usepackage{fontspec}
\usepackage[spanish]{babel}
\usepackage{tikz}
% Librerías TikZ para los diagramas de las guías (bloque `recursos:`):
%  · babel  → babel en español vuelve activos caracteres como «>», y TikZ los usa
%             en sus opciones (>=stealth, ->). Sin esta librería, cualquier
%             diagrama con flechas revienta con «\language@active@arg> has an
%             extra }». Debe ir ANTES de que se dibuje nada.
%  · positioning        → sintaxis «below left=2pt and 3pt».
%  · decorations.pathreplacing → llaves ({brace}) para señalar rangos.
\usetikzlibrary{babel,positioning,decorations.pathreplacing}
\usepackage{graphicx}
% Circuitos: el trabajo de electrónica se hace hoy con micro:bit y sus sensores
% integrados (MakeCode, sin cableado), así que no se necesitan esquemáticos.
% Si algún día entran montajes con protoboard, basta descomentar esta línea:
% los diagramas .tex/.tikz declarados en `recursos:` ya se incrustan solos.
% \usepackage{circuitikz}
\usepackage[most]{tcolorbox}
\usepackage{tabularx}
\usepackage{array}
\usepackage{fancyhdr}
\usepackage{titlesec}
\usepackage{ragged2e}
\usepackage{enumitem}
\usepackage{microtype}

% Helvetica Neue no trae la flecha U+2192, asi que cada «→» escrito literalmente
% en el YAML se compilaba a NADA, en silencio (462 flechas en 61 guias: rutas del
% tipo «paleta → tipografia → lockup» salian rotas). Mapeamos el caracter a su
% version matematica, que si tiene glifo. Asi el autor puede seguir escribiendo
% «→» comodamente en el YAML.
\usepackage{newunicodechar}
\newunicodechar{→}{\ensuremath{\rightarrow}}

\setmainfont{Helvetica Neue}
\setsansfont{Helvetica Neue}
\setlength{\parindent}{0pt}
\setlength{\parskip}{6pt}
\hyphenpenalty=200
\exhyphenpenalty=200
\pretolerance=400
\tolerance=2000
\setlength{\emergencystretch}{1em}
\renewcommand{\arraystretch}{1.25}
\setlist{leftmargin=5mm,itemsep=1mm,topsep=1mm}
\newcolumntype{Y}{>{\RaggedRight\arraybackslash}X}

\definecolor{milcMagenta}{HTML}{E6007E}
\definecolor{milcVino}{HTML}{5A0038}
\definecolor{milcTurquesa}{HTML}{0093A5}
\definecolor{milcVerde}{HTML}{58A923}
\definecolor{milcMostaza}{HTML}{E5B400}
\definecolor{milcOcre}{HTML}{B86600}
\definecolor{milcNegro}{HTML}{241816}
\definecolor{milcGris}{HTML}{F7F7F4}
\definecolor{milcLinea}{HTML}{E8E2DD}
\definecolor{milcGradePrimary}{HTML}{EA580C}
\definecolor{milcGradeDark}{HTML}{9A3412}
\definecolor{milcGradeSoft}{HTML}{FFEDD5}
\definecolor{milcGradeAccent}{HTML}{CFE7FF}
\definecolor{milcGradeText}{HTML}{FFFFFF}
\color{milcNegro}

\pagestyle{fancy}
\fancyhf{}
\lhead{\small Semillero de Investigación · Robótica y sistemas embebidos}
\rhead{\small Módulo 1 · Escucha · MILC v3}
\cfoot{\small\thepage}
\renewcommand{\headrulewidth}{0pt}

\newtcolorbox{titlebox}[1]{
  enhanced,colback=#1,colframe=#1,arc=7mm,boxrule=0pt,
  left=7mm,right=7mm,top=3mm,bottom=3mm,
  before skip=8pt,after skip=8pt,
  fontupper=\sffamily\bfseries\Large\color{white}
}

\newtcolorbox{softbox}[3]{
  enhanced,colback=#2,colframe=#1,borderline west={4pt}{0pt}{#1},
  arc=2mm,boxrule=.4pt,left=4mm,right=4mm,top=3mm,bottom=3mm,
  before skip=6pt,after skip=8pt,title={#3},coltitle=white,
  colbacktitle=#1,fonttitle=\sffamily\bfseries,fontupper=\RaggedRight,
  attach boxed title to top left={xshift=3mm,yshift=-2mm},
  boxed title style={arc=2mm, boxrule=0pt}
}

\newtcolorbox{drawbox}[2]{
  enhanced,colback=white,colframe=#1,arc=4mm,boxrule=1pt,
  left=5mm,right=5mm,top=4mm,bottom=4mm,before skip=8pt,after skip=8pt,
  title={#2},coltitle=white,colbacktitle=#1,fonttitle=\sffamily\bfseries,
  fontupper=\RaggedRight,
  attach boxed title to top left={xshift=4mm,yshift=-2mm},
  boxed title style={arc=3mm, boxrule=0pt}
}

\newtcolorbox{infoband}[1]{
  enhanced,colback=#1!8,colframe=#1!50,arc=2mm,boxrule=.5pt,
  left=4mm,right=4mm,top=2mm,bottom=2mm,before skip=4pt,after skip=4pt,
  fontupper=\small\RaggedRight
}

% Puente narrativo entre secciones (contrato editorial Conéctate).
% Da continuidad: "Acabas de X. Ahora vas a Y porque Z."
% Visualmente: banda izquierda en color del grado, texto en itálica pequeña.
\newtcolorbox{puentebox}{
  enhanced,colback=milcGradeSoft,colframe=milcGradeDark,
  borderline west={3pt}{0pt}{milcGradeDark},
  arc=0mm,boxrule=0pt,left=5mm,right=4mm,top=2.5mm,bottom=2.5mm,
  before skip=4pt,after skip=8pt,
  fontupper=\itshape\small\color{milcNegro}\RaggedRight
}
% La flecha va en modo matematico a proposito: Helvetica Neue NO tiene el glifo
% U+2192, asi que \textrightarrow se compilaba a NADA («Missing character: There
% is no → in font Helvetica Neue») y el puente salia sin flecha. En modo
% matematico la flecha viene de la fuente matematica y si se dibuja.
\newcommand{\puente}[1]{%
  \begin{puentebox}%
    \textcolor{milcGradeDark}{$\rightarrow$}\hspace{2mm}#1%
  \end{puentebox}%
}

\newcommand{\linead}[1]{\noindent\color{milcLinea}\rule{#1}{0.5pt}\color{milcNegro}}
\newcommand{\respuesta}[1]{\par\vspace{2mm}\foreach \n in {1,...,#1}{\noindent\color{milcLinea}\rule{\linewidth}{0.5pt}\par\vspace{5mm}}\color{milcNegro}}
\newcommand{\checkbox}{\textcolor{milcGradeDark}{\fbox{\phantom{x}}}\hspace{5pt}}

\newcommand{\phasecard}[4]{
\begin{tcolorbox}[enhanced,colback=#3,colframe=#2,arc=3mm,boxrule=.5pt,height=3.05cm,valign=center,halign=center,left=1.5mm,right=1.5mm,top=2mm,bottom=2mm]
{\sffamily\bfseries\fontsize{22}{24}\selectfont\textcolor{#2}{#1}}\par
{\sffamily\bfseries\small #4}
\end{tcolorbox}
}

% ── Recursos visuales de la guía ──────────────────────────────────────
% Declarados en el bloque `recursos:` del YAML y emitidos por el builder.
% \guiaFigura   → imagen rasterizada o PDF (foto, carta celeste, montaje).
% \guiaDiagrama → fuente TikZ/circuitikz (.tex) incrustada como vector:
%                 es la vía para circuitos y esquemáticos editables.
% #1 = ruta relativa al .tex · #2 = pie (puede ir vacío)
\newcommand{\guiaFigura}[2]{%
\begin{center}
\includegraphics[width=0.88\linewidth,height=0.38\textheight,keepaspectratio]{#1}\\[1.5mm]
{\footnotesize\itshape\textcolor{milcNegro!75}{#2}}
\end{center}
}

\newcommand{\guiaDiagrama}[2]{%
\begin{center}
\input{#1}\\[1.5mm]
{\footnotesize\itshape\textcolor{milcNegro!75}{#2}}
\end{center}
}

\begin{document}
\thispagestyle{empty}
\begin{tikzpicture}[remember picture,overlay]
  \fill[white] (current page.south west) rectangle (current page.north east);
  \path[fill=milcGradePrimary,rounded corners=16mm]
    ([xshift=.55cm,yshift=-.55cm]current page.north west) rectangle
    ([xshift=-.55cm,yshift=3.65cm]current page.south east);

  \node[anchor=north west,text=milcGradeText] at ([xshift=2.0cm,yshift=-1.85cm]current page.north west)
    {\sffamily\bfseries\fontsize{14}{16}\selectfont MÓDULO};
  \node[anchor=north west,text=milcGradeText,opacity=.82] at ([xshift=2.0cm,yshift=-2.75cm]current page.north west)
    {\sffamily\bfseries\fontsize{9}{11}\selectfont SEMILLERO DE INVESTIGACIÓN · CONECTATE};

  \begin{scope}[shift={([xshift=-3.05cm,yshift=-2.55cm]current page.north east)}]
    \fill[white,opacity=.80] (-.62,-.52) rectangle (.62,.52);
    \draw[milcGradeText,line width=1pt] (-.62,-.52) rectangle (.62,.52);
    \fill[milcGradeDark] (-.42,-.38) rectangle (-.22,.06);
    \fill[milcGradeAccent] (-.10,-.38) rectangle (.10,.24);
    \fill[milcGradePrimary!60!black] (.22,-.38) rectangle (.42,.38);
  \end{scope}

  \node[anchor=west,text=milcGradeText] at ([xshift=1.9cm,yshift=-12.85cm]current page.north west)
    {\sffamily\bfseries\fontsize{124}{124}\selectfont 1};
  % El circulito de grado (°) solo aplica a las guias de grado («11°»); en
  % Bebras y Semillero el numero grande es un momento/modulo, no un grado.
  

  \node[anchor=west,text=milcGradeText,text width=8.7cm,align=left] at ([xshift=10.35cm,yshift=-11.6cm]current page.north west)
    {
     {\sffamily\bfseries\fontsize{19}{22}\selectfont Cosas que podrían\\reaccionar ---\\del sapo que anuncia\\la lluvia\\al sensor que percibe}\\[4mm]
     {\sffamily\bfseries\fontsize{23}{26}\selectfont Robótica y sistemas embebidos}\\[2mm]
     {\sffamily\fontsize{13}{16}\selectfont Módulo 1 · Fase MILC: Escucha · 3 h de clase}\\[5mm]
     {\sffamily\bfseries\fontsize{11}{13}\selectfont Institución Educativa Sor María Juliana}\\[1mm]
     {\sffamily\fontsize{10}{12}\selectfont Autor: Dr. Álvaro Cárdenas Orozco}
    };

  \path[fill=milcGradeSoft,rounded corners=7mm]
    ([xshift=1.35cm,yshift=.85cm]current page.south west) rectangle
    ([xshift=-1.35cm,yshift=4.25cm]current page.south east);
  \draw[milcGradeDark,line width=.65pt,rounded corners=7mm]
    ([xshift=1.35cm,yshift=.85cm]current page.south west) rectangle
    ([xshift=-1.35cm,yshift=4.25cm]current page.south east);
  \node[anchor=center,text=milcNegro,text width=17.0cm,align=center] at ([yshift=2.55cm]current page.south)
    {\sffamily\fontsize{10}{12}\selectfont Metodología MILC · Apertura ancestral · Pensamiento computacional · Triángulo Dussel-Estoicismo-Floridi\\
    \textcolor{milcGradeDark}{\bfseries Producto:} Tu \textbf{mapa de necesidades y sensores}: una lista de \textbf{necesidades reales} de tu entorno (casa, escuela, vereda) que un aparato \emph{que percibe y actúa} podría atender, y al lado, \textbf{con cuál sensor del micro:bit} se alcanzaría cada una (luz, temperatura, movimiento, botones, brújula). Más \textbf{una necesidad elegida y afinada}: qué debe \textbf{percibir} el aparato y qué debe \textbf{hacer} en respuesta ---la semilla del prototipo que construirás en el módulo 3.};
\end{tikzpicture}
\mbox{}
\newpage

\section*{Apertura: Cosas que podrían reaccionar --- del sapo que anuncia la lluvia al sensor que percibe}

\begin{softbox}{milcGradeDark}{milcGradeSoft}{Saber ancestral}
\textbf{Para el pueblo Nasa ---del Cauca y el sur del Valle--- la naturaleza avisa, y hay que saber escucharla.} El canto del \textbf{sapo} y de la \textbf{rana}, el paso del \textbf{colibrí}, ciertas \textbf{culebras}: los Nasa los reconocen como seres ligados a la energía de la \textbf{lluvia}; su aparición o su canto anuncia que el agua viene (Ulloa, 2011). No es superstición: es \textbf{lectura atenta de bioindicadores} ---seres vivos que sienten el cambio del ambiente \emph{antes} que nosotros y reaccionan, y que el pueblo aprendió a leer para decidir cuándo sembrar o cuándo resguardarse. Míralo de cerca: cada uno de esos seres es, en el fondo, un \textbf{sensor vivo}. Percibe un cambio del entorno (la humedad, la presión que trae la lluvia) y \textbf{reacciona} (canta, se mueve, aparece); y la comunidad, al \textbf{leer} esa reacción, actúa a tiempo. Un aparato robótico hace exactamente lo mismo con otros sensores: \textbf{percibe una señal del entorno y responde}. Pero antes de construir nada hay que aprender lo más difícil, lo que los Nasa dominan: \textbf{escuchar} el entorno, reconocer qué está pasando y qué pide respuesta. La máquina que percibe y actúa no empieza en un cable: empieza en una \textbf{pregunta}.
\end{softbox}

\begin{softbox}{milcTurquesa}{milcGris}{Saber contemporáneo}
Un \textbf{sistema embebido} es un computador escondido dentro de un objeto que \textbf{percibe} su entorno y \textbf{actúa} sobre él ---el que enciende la luz cuando entras, el que apaga la nevera al enfriar, el que abre la puerta al detectar tu mano. Todos tienen la misma anatomía de tres partes: \textbf{sensor} (percibe: luz, calor, movimiento) → \textbf{lógica} (decide con reglas) → \textbf{actuador} (hace: prende, suena, muestra, mueve). Y aquí está la buena noticia del \textbf{micro:bit}: \textbf{ya trae los sensores adentro}, sin cablear nada. Su matriz de 25 LED también \textbf{mide la luz}; incluye \textbf{sensor de temperatura}, \textbf{acelerómetro} (movimiento, agitar, caída, inclinación), \textbf{brújula} (dirección) y dos \textbf{botones}. Cada sensor es una forma de \emph{escuchar} el mundo, como el sapo escucha la lluvia. Pero un buen investigador no empieza preguntando \emph{``¿qué puedo hacer con este sensor?''} ---eso es enamorarse del juguete. Empieza al revés, como el Nasa: \textbf{¿qué de mi entorno necesita ser atendido?}, y solo entonces \emph{¿qué sensor lo alcanza?}. La técnica que sirve nace de una necesidad real, no de un aparato en busca de para qué.
\end{softbox}

\begin{softbox}{milcMostaza}{milcGris}{Pregunta-puente}
\textit{El Nasa no leía al sapo por curiosidad: lo leía para \textbf{decidir a tiempo} qué hacer con su siembra; \textbf{¿qué necesidad de tu casa, tu escuela o tu vereda podría «escuchar» un aparato con sensores}\ldots y qué haría en respuesta, para que sirva de verdad?}
\end{softbox}

\begin{softbox}{milcVerde}{milcGris}{Saber y saber hacer}
Al terminar la sesión: (1) podrás \textbf{identificar} necesidades reales de tu entorno que un sistema que percibe y actúa podría atender; (2) sabrás \textbf{explicar} la anatomía sensor → lógica → actuador y qué percibe cada sensor del micro:bit; (3) podrás \textbf{aplicar} ese mapa relacionando cada necesidad con el sensor que la alcanza; (4) habrás \textbf{definido} una necesidad elegida: qué debe percibir tu aparato y qué debe hacer.
\end{softbox}



\small
\begin{tabularx}{\textwidth}{>{\bfseries}p{3.0cm}Y}
\rowcolor{milcGradePrimary!12}Autor & Dr. Álvaro Cárdenas Orozco\\
DBA & Formula preguntas investigables, produce evidencia con método y comunica hallazgos reconociendo sus límites (investigación estudiantil STEM, transversal 6°--11°)\\
Referentes & micro:bit · MakeCode · Continuidad con la unidad de 8° P2 (guías 8-2-1 a 8-2-10)\\
Producto final & Tu \textbf{mapa de necesidades y sensores}: una lista de \textbf{necesidades reales} de tu entorno (casa, escuela, vereda) que un aparato \emph{que percibe y actúa} podría atender, y al lado, \textbf{con cuál sensor del micro:bit} se alcanzaría cada una (luz, temperatura, movimiento, botones, brújula). Más \textbf{una necesidad elegida y afinada}: qué debe \textbf{percibir} el aparato y qué debe \textbf{hacer} en respuesta ---la semilla del prototipo que construirás en el módulo 3.\\
\end{tabularx}
\normalsize

\puente{Arranca la línea de Robótica. Hoy no construyes nada todavía: aprendes lo que el Nasa domina y el ingeniero necesita ---\textbf{escuchar el entorno} para descubrir qué necesidad podría atender un aparato que percibe y actúa.}

\begin{titlebox}{milcGradeDark}
Ruta MILC: así vamos a aprender
\end{titlebox}
\begin{tabularx}{\textwidth}{>{\centering\arraybackslash}X>{\centering\arraybackslash}X>{\centering\arraybackslash}X>{\centering\arraybackslash}X}
\phasecard{1}{milcVerde}{milcGris}{Escucha\\[-1mm]\small Observo, escucho y pregunto.} &
\phasecard{2}{milcTurquesa}{milcGris}{Sistematización\\[-1mm]\small Pensamiento computacional.} &
\phasecard{3}{milcMagenta}{milcGris}{Praxis\\[-1mm]\small Construyo y aplico.} &
\phasecard{4}{milcVino}{milcGris}{Evaluación\\[-1mm]\small Reflexión liberadora.}
\end{tabularx}


\puente{Antes de pensar en aparatos, mira tu entorno con atención. Como el que lee al sapo: las necesidades ya están ahí, esperando que alguien las note.}

\begin{titlebox}{milcVerde}
Fase 1 · Escucha
\end{titlebox}

\begin{softbox}{milcVerde}{milcGris}{Escena para escuchar}
\textbf{Actividad 1 · IDENTIFICA --- Detective de necesidades} (40 min · individual + grupo). \textbf{Momento A (10 min) --- El caso del sapo.} El profe cuenta cómo los Nasa leen al sapo, la rana y el colibrí como señales de lluvia, y pregunta: \emph{¿qué es un «sensor vivo» y cómo la comunidad actúa al leerlo?}. \textbf{Momento B (15 min) --- Safari de necesidades.} Cada quien recorre con la mente su casa, su escuela y su vereda buscando \textbf{molestias, riesgos o desperdicios} que un aparato \emph{que percibe y actúa} podría atender: \emph{``la puerta del corral queda abierta y no nos damos cuenta''}, \emph{``el tanque se rebosa''}, \emph{``la abuela no oye el timbre''}, \emph{``dejamos la luz prendida''}. Anota \textbf{mínimo 6}. \textbf{Momento C (15 min) --- Necesidad vs capricho.} En grupo, separen las que \textbf{le sirven a alguien de verdad} de las que son solo \emph{``sería chévere''}. \textbf{En el cuaderno:} tus 6+ necesidades + la marca de cuáles son reales y a quién le sirven.
\end{softbox}

\checkbox Puedo explicar qué es un «sensor vivo» con el ejemplo del sapo Nasa.\\
\checkbox Anoté al menos 6 necesidades reales de mi entorno.\\
\checkbox Separé las necesidades que le sirven a alguien de los caprichos.

\begin{infoband}{milcVerde}


\textbf{Lo que va al cuaderno (Actividad 1 · IDENTIFICA).} Título: «Actividad 1 --- Detective de necesidades». \textbf{Formato:} lista de 6+ necesidades del entorno + a quién le sirve cada una + la marca «real / capricho». \textbf{Extensión:} 10 líneas. \textbf{Sabes que terminaste cuando} puedes señalar \textbf{una necesidad que le importa a una persona concreta} (no a ti por diversión).
\end{infoband}

\puente{Ya tienes necesidades detectadas. Ahora aprende con qué se perciben: la anatomía sensor → lógica → actuador y los sensores que el micro:bit trae adentro.}

\begin{titlebox}{milcTurquesa}
Fase 2 · Sistematización con pensamiento computacional
\end{titlebox}

\begin{softbox}{milcTurquesa}{milcGris}{Cuatro pilares aplicados al tema}
\textbf{Actividad 2 · EXPLICA --- Con qué se escucha el mundo} (65 min · grupo + individual). Ya tienes necesidades. Ahora aprende la \textbf{anatomía} de un aparato que percibe y actúa, y \textbf{qué percibe cada sensor} del micro:bit ---sus «sentidos» integrados.
\end{softbox}

\small
\begin{tabularx}{\textwidth}{>{\bfseries\RaggedRight\arraybackslash}p{3.6cm}Y}
\rowcolor{milcTurquesa!90}\color{white}Pilar & \color{white}\bfseries Aplicación al tema\\
1. Descomposición & \textbf{Pilar 1 --- Sensor → lógica → actuador: la anatomía de todo.} Desde una puerta automática hasta un satélite, todo sistema que reacciona tiene \textbf{tres partes}. El \textbf{sensor} \emph{percibe} una señal del entorno (hay luz, hace calor, algo se movió) ---es el oído, el ojo, el tacto del aparato. La \textbf{lógica} \emph{decide} qué hacer con esa señal, comparándola con una regla (\emph{si está oscuro, entonces\ldots}). El \textbf{actuador} \emph{hace} algo visible: prende un LED, suena, muestra un ícono, mueve un motor. \textbf{Quítale cualquiera de las tres y deja de funcionar:} un sensor sin lógica no decide, una lógica sin actuador no hace, un actuador sin sensor no sabe cuándo. El sapo Nasa tenía las tres: piel que siente (sensor), instinto que decide (lógica), canto que avisa (actuador).\\
2. Reconocimiento de patrones & \textbf{Pilar 2 --- Los sentidos del micro:bit vienen adentro.} La gran ventaja para investigar sin presupuesto: el micro:bit \textbf{ya trae sensores}, sin cablear ni comprar nada. \textbf{Luz:} su matriz de 25 LED, además de mostrar, \emph{mide} la luz que le llega. \textbf{Temperatura:} siente el calor del ambiente. \textbf{Acelerómetro:} detecta \emph{movimiento, agitar, caída e inclinación} ---sabe si lo mueves, lo volteas o se cae. \textbf{Brújula:} conoce la \emph{dirección} (norte, campo magnético). \textbf{Botones A y B:} reciben la orden de una persona. \textbf{Pines:} permiten el tacto y conectar sensores externos si algún día hicieran falta. \emph{Cada uno es una forma de escuchar el mundo}, como el pueblo Nasa lee distintas señales: el sapo para la lluvia, el viento para la tormenta, las estrellas para el tiempo.\\
3. Abstracción & \textbf{Pilar 3 --- Empieza por la necesidad, no por el sensor.} Este es el error que separa al ingeniero del que juega. El principiante dice: \emph{``tengo un acelerómetro, ¿qué hago con él?''} ---y termina con un aparato que no le sirve a nadie. El investigador hace lo contrario, como el Nasa: parte de una \textbf{necesidad real} (\emph{``hay que saber si la puerta del corral quedó abierta''}) y \emph{después} busca qué sensor la alcanza (\emph{``el acelerómetro puede sentir si la puerta se movió''}). \textbf{La necesidad manda; el sensor obedece.} Un aparato nace bien cuando responde a una pregunta del entorno, y nace mal cuando es un juguete buscando para qué.\\
4. Algoritmo conceptual & \textbf{Pilar 4 --- Cruzar necesidad y sensor: el mapa.} La herramienta de hoy es una \textbf{tabla de dos columnas}: a la izquierda, la necesidad; a la derecha, \textbf{qué sensor la percibiría} y \textbf{qué haría el actuador}. Ejemplos: \emph{``dejamos la luz prendida de día''} → sensor de \textbf{luz} + actuador que \emph{avisa} cuando hay luz de sol y una lámpara encendida. \emph{``el tanque se rebosa''} → sensor (con pines) + actuador que \emph{suena}. \emph{``la puerta del corral quedó abierta''} → \textbf{acelerómetro} + actuador que \emph{muestra alerta}. Algunas necesidades \textbf{no} las alcanza el micro:bit solo ---y reconocerlo también es parte del mapa. \emph{Ese cruce honesto entre lo que se necesita y lo que el aparato puede es el corazón del diseño.}\\
\end{tabularx}
\normalsize

\begin{drawbox}{milcTurquesa}{Sensor → lógica → actuador + sentidos del micro:bit}
\textbf{Anatomía:} sensor (percibe) → lógica (decide con una regla) → actuador (hace). Quita una y deja de funcionar. \\[1mm]
\textbf{Sensores del micro:bit (sin cablear):} luz (matriz LED) · temperatura · acelerómetro (movimiento/caída/inclinación) · brújula (dirección) · botones A/B · pines (tacto). \\[1mm]
\textbf{Regla de oro:} primero la necesidad, después el sensor. La necesidad manda. \\[1mm]
\textbf{Mapa:} tabla necesidad → sensor que la percibe → actuador que responde.
\end{drawbox}

\begin{softbox}{milcMostaza}{milcGris}{Errores comunes y su corrección}
\textbf{Errores al empezar un proyecto de robótica (y cómo evitarlos).} (1) \emph{Enamorarse del sensor} --- «tengo el acelerómetro, algo haré»; \textbf{parte de la necesidad}, no del aparato. (2) \emph{Olvidar el actuador} --- percibir sin hacer nada no es un sistema; define \textbf{qué acción visible} responde. (3) \emph{Necesidad-capricho} --- «que titile bonito» no le sirve a nadie; busca una necesidad con \textbf{dueño}. (4) \emph{Pedirle al micro:bit lo que no puede} --- no todo sensor existe integrado; reconocer el \textbf{límite} es parte del mapa. (5) \emph{Confundir sensor con actuador} --- el sensor \emph{entra} información, el actuador \emph{saca} acción; no son lo mismo. (6) \emph{Saltar a los bloques} --- todavía no se programa; hoy se \textbf{escucha y se mapea}. El código llega en el módulo 3.
\end{softbox}

\begin{infoband}{milcTurquesa}
\guiaDiagrama{assets/robotica-1/sensores-microbit.tex}{Los sentidos del micro:bit vienen adentro, sin cablear: luz, temperatura, movimiento, dirección y tacto. Foto: Gareth Halfacree, vía Wikimedia Commons (CC BY-SA 2.0).}

\textbf{Lo que va al cuaderno (Actividad 2 · EXPLICA).} Título: «Actividad 2 --- Con qué se escucha el mundo». \textbf{Formato:} 3 bloques. Bloque A: el esquema sensor → lógica → actuador con un ejemplo cotidiano. Bloque B: la lista de sensores del micro:bit y qué percibe cada uno. Bloque C: 2 necesidades tuyas ya cruzadas con su sensor y su actuador. \textbf{Extensión:} 1 hoja. \textbf{Sabes que terminaste cuando} podrías explicarle a alguien \textbf{por qué se empieza por la necesidad y no por el sensor}.
\end{infoband}

\puente{Ya sabes qué percibe cada sensor. Es hora de \textbf{cruzar} tus necesidades con los sensores y elegir una: qué debe percibir tu aparato y qué debe hacer.}

\begin{titlebox}{milcMagenta}
Fase 3 · Praxis con pensamiento computacional
\end{titlebox}

\begin{softbox}{milcMagenta}{milcGris}{Construyendo el producto con los cuatro pilares}
\textbf{Actividad 3 · APLICA --- Arma tu mapa y elige tu necesidad} (45 min · individual + parejas). Hora de cruzar todo. Vas a convertir tus necesidades en un \textbf{mapa de necesidad → sensor → actuador}, y a \textbf{elegir una} para afinarla: la semilla del prototipo del módulo 3.
\end{softbox}

\small
\begin{tabularx}{\textwidth}{>{\bfseries\RaggedRight\arraybackslash}p{3.6cm}Y}
\rowcolor{milcMagenta!90}\color{white}Pilar & \color{white}\bfseries Aplicación al producto\\
Descomposición & \textbf{Paso 1 --- Llena el mapa (15 min).} Toma tus 6+ necesidades de la Actividad 1 y arma la tabla de \textbf{tres columnas}: \emph{Necesidad · Qué sensor la percibe · Qué haría el actuador}. Sé concreto: no «un sensor», sino \emph{cuál} (luz, temperatura, acelerómetro, botón). Si una necesidad \textbf{no} la alcanza el micro:bit solo, escríbela igual y márcala \emph{``necesita sensor externo''} ---ese reconocimiento vale.\\
Patrones & \textbf{Paso 2 --- Prueba el «test del sapo» (10 min).} Por cada fila, pregúntate como el Nasa: \emph{¿esto responde a una necesidad real, con dueño, o es un capricho bonito?}. Tacha las de capricho. De las que quedan, marca \textbf{la que más le importa a una persona concreta} (tu abuela, tu vereda, tu curso). Esa es tu candidata.\\
Abstracción & \textbf{Paso 3 --- Afina tu necesidad elegida (15 min).} Escríbela en el \textbf{formato del ingeniero}: \emph{``El aparato debe PERCIBIR \_\_\_\_ (qué señal, con qué sensor) y, cuando \_\_\_\_ (la regla), debe HACER \_\_\_\_ (qué acción visible)''}. Ejemplo: \emph{``debe percibir la luz con el sensor de luz y, cuando haya sol Y la lámpara siga encendida, debe mostrar una alerta''}. \textbf{Esa frase es el plano de tu prototipo:} en el módulo 2 la volverás pseudocódigo, en el 3 la construirás.\\
Algoritmo de armado & \textbf{Paso 4 --- Dibuja el sistema (5 min).} Haz un dibujo simple de tres cajas ---\textbf{sensor → lógica → actuador}--- con tu necesidad: qué entra por el sensor, qué regla decide, qué sale por el actuador. Un dibujo que cualquiera entienda de un vistazo. \emph{Si no puedes dibujarlo, todavía no está claro; afínalo hasta que sí.}\\
\end{tabularx}
\normalsize

\begin{drawbox}{milcMagenta}{Antes de cerrar tu mapa}
\checkbox Mi tabla cruza cada necesidad con un sensor concreto y un actuador. \\[1mm]
\checkbox Marqué qué necesidades no alcanza el micro:bit solo (y no las borré). \\[1mm]
\checkbox Pasé el «test del sapo»: mi necesidad elegida le sirve a una persona concreta. \\[1mm]
\checkbox Escribí mi necesidad en formato PERCIBIR → REGLA → HACER. \\[1mm]
\checkbox Dibujé el sistema en tres cajas (sensor → lógica → actuador). \\[1mm]
\checkbox Sé cuál es la semilla de mi prototipo para el módulo 3.
\end{drawbox}

\begin{softbox}{milcMagenta}{milcGris}{Plantilla de trabajo}
\textbf{Plantilla del mapa.} \textit{Tabla:} Necesidad · Sensor que la percibe · Actuador que responde · ¿A quién le sirve?. \\[1mm]
\textit{Necesidad elegida (formato ingeniero):} debe PERCIBIR \_\_\_\_ (sensor \_\_\_\_), y cuando \_\_\_\_ (regla), debe HACER \_\_\_\_ (acción). \\[1mm]
\textit{Dibujo:} [ sensor ] → [ lógica: la regla ] → [ actuador ].
\end{softbox}

\begin{infoband}{milcMagenta}


\textbf{Lo que va al cuaderno (Actividad 3 · APLICA).} Título: «Actividad 3 --- Mi mapa de necesidad y sensor». \textbf{Formato:} la tabla necesidad → sensor → actuador + el «test del sapo» aplicado + la necesidad elegida en formato PERCIBIR/REGLA/HACER + el dibujo de tres cajas. \textbf{Extensión:} 1 hoja. \textbf{Sabes que terminaste cuando} tu necesidad elegida dice claramente qué percibe, con qué regla decide y qué hace.
\end{infoband}

\puente{El producto no es un aparato todavía: es un \textbf{mapa de necesidades y sensores} y una necesidad afinada. Sin eso, construir es enamorarse del juguete.}

\begin{titlebox}{milcMostaza}
Producto de la sesión
\end{titlebox}
\textbf{Mapa de necesidades y sensores + una necesidad afinada (semilla del prototipo)}.
\begin{softbox}{milcMostaza}{milcGris}{Criterios de calidad}
(1) La tabla cruza cada necesidad con un sensor concreto del micro:bit y un actuador. (2) Están marcadas las necesidades que el micro:bit no alcanza solo (no se ocultan). (3) La necesidad elegida le sirve a una persona concreta (pasó el «test del sapo»). (4) La necesidad está escrita en formato PERCIBIR → REGLA → HACER. (5) Hay un dibujo del sistema en tres cajas (sensor → lógica → actuador).
\end{softbox}

\puente{Antes de cerrar, prueba tu necesidad elegida: ¿de verdad le sirve a alguien, o solo es «lo que se puede hacer con el sensor»?}

\begin{titlebox}{milcVino}
Fase 4 · Evaluación liberadora — cinco dimensiones
\end{titlebox}

\begin{softbox}{milcVino}{milcGris}{Sentido de la evaluación}
Evaluar no es solo revisar una nota. Es reconocer qué aprendí, qué cambió en mí, qué emociones manejé, cómo me relaciono con la ciudad y la comunidad, y qué le dejaré a quien viene detrás.
\end{softbox}

\textbf{Mi reflexión liberadora en cinco dimensiones.} Completa cada frase iniciada:

\small
\begin{tabularx}{\textwidth}{>{\bfseries\RaggedRight\arraybackslash}p{3.4cm}Y>{\RaggedRight\arraybackslash}p{4.6cm}}
\rowcolor{milcVino}\color{white}Dimensión & \color{white}\bfseries Mi respuesta (frame starter) & \color{white}\bfseries Algo que descubrí\\
1. Personal & Lo que más cambió en mí fue \linead{6.5cm} & \linead{4.2cm}\\
2. Emocional & La emoción más fuerte fue \linead{2.5cm} y la regulé \linead{3cm} & \linead{4.2cm}\\
3. Ciudadana & En democracia esto importa porque \linead{6cm} & \linead{4.2cm}\\
4. Local (Cartago) & En mi barrio sería útil para \linead{6.5cm} & \linead{4.2cm}\\
5. Intergeneracional & A mi abuela/abuelo le diría \linead{6.5cm} & \linead{4.2cm}\\
\end{tabularx}
\normalsize

\begin{infoband}{milcVino}
\textbf{Para la próxima sustentación.} Vuelve a esta tabla en seis meses. Verás que las respuestas cambian — no porque lo aprendido cambie, sino porque tú cambias. La evaluación liberadora es una conversación contigo a través del tiempo.
\end{infoband}


\puente{Cerramos con 3 voces que te ayudan a entender por qué la buena técnica empieza escuchando una necesidad real, no luciendo un aparato.}

\begin{titlebox}{milcGradeDark}
Triángulo de pensamiento · cierre filosófico
\end{titlebox}

\begin{softbox}{milcVino}{milcGris}{Dussel · lente del nosotros}
\textit{``El otro se revela realmente como otro\ldots como el pobre, el oprimido; el que, a la vera del camino, fuera del sistema, muestra su rostro sufriente y sin embargo desafiante: «¡Tengo hambre!, ¡tengo derecho a comer!»''}

Dussel dice que el otro se revela reclamando una necesidad. Un aparato que sirve nace de escuchar ese reclamo: la abuela que no oye el timbre, el corral que queda abierto, el tanque que se rebosa. \textbf{La técnica que empieza en una necesidad real ---la del otro, no la de mi lucimiento--- es técnica al servicio de la vida.} Enamorarse del sensor es hablar solo; escuchar la necesidad es dejar que el otro hable primero.

\textbf{¿Mi idea responde al reclamo real de alguien, o a mis ganas de lucir un aparato?}
\end{softbox}
\linead{\linewidth}\\[1mm]
\linead{\linewidth}

\begin{softbox}{milcMostaza}{milcGris}{Estoicismo · Marco Aurelio · Meditaciones VII, 47 (c. 175 d.C.) · cuidado interior}
\textit{``Contempla las evoluciones de los astros y piensa al mismo tiempo que evolucionas con ellos; y reflexiona continuamente cómo los elementos cambian unos en otros.''}

Marco Aurelio pide mirar el mundo \textbf{poniéndote adentro}, como parte de lo que observas. Antes de construir cualquier aparato, la fase de \emph{escucha} pide justo eso: mirar tu entorno con atención ---la casa, la escuela, la vereda--- hasta \textbf{ver las necesidades que estaban ahí y nadie notaba}. El Nasa que lee al sapo no mira desde afuera: está adentro del mismo mundo que el sapo anuncia. \textbf{Percibir bien es el primer acto técnico, y es un acto de atención.}

\textbf{¿Miro mi entorno con la atención suficiente para notar lo que otros pasan por alto?}
\end{softbox}
\linead{\linewidth}\\[1mm]
\linead{\linewidth}

\begin{softbox}{milcTurquesa}{milcGris}{Floridi · lente de la infoesfera}
\textit{``Somos organismos informacionales ---inforgs--- que compartimos un entorno global hecho, en última instancia, de información: la infoesfera.''}

Floridi dice que vivimos rodeados de información. Un sensor es, precisamente, una \textbf{puerta} por la que un pedazo del mundo ---la luz, el calor, el movimiento--- se vuelve \emph{información} que un aparato puede usar para decidir. Percibir es transformar mundo en dato. \textbf{Cuando eliges qué sensor poner, decides qué parte del entorno vas a dejar entrar} a la máquina; por eso empezar por la necesidad, y no por el sensor, decide qué mundo va a «escuchar» tu aparato.

\textbf{¿Qué parte de mi entorno estoy dejando que mi aparato perciba, y por qué esa y no otra?}
\end{softbox}
\linead{\linewidth}\\[1mm]
\linead{\linewidth}

\begin{titlebox}{milcVino}
Compromiso final
\end{titlebox}

\small
\begin{tabularx}{\textwidth}{>{\RaggedRight\arraybackslash}p{3.0cm}YYY}
\rowcolor{milcVino}\color{white}\bfseries Criterio & \color{white}\bfseries Lo logré & \color{white}\bfseries En proceso & \color{white}\bfseries Necesito apoyo\\
Comprensión del tema & Explico ideas centrales con mis palabras. & Reconozco ideas pero falta precisión. & Necesito repasar conceptos base.\\
Pensamiento computacional & Apliqué los 4 pilares al diseñar mi producto. & Apliqué algunos pilares. & Necesito releer la fase 2.\\
Aplicación y producto & Mi producto cumple los criterios y se puede revisar. & Producto funcional con ajustes pendientes. & Aún en construcción.\\
Reflexión 5 dimensiones & Respondí cada dimensión con honestidad. & Respondí algunas con detalle. & Mis respuestas son muy generales.\\
Triángulo de pensamiento & Conecté las 3 voces con mi trabajo real. & Conecté al menos una. & No relaciono las voces con mi proceso.\\
\end{tabularx}
\normalsize

\textbf{Hoy entendí que:}\\[1mm]
\linead{\linewidth}\\[1mm]
\linead{\linewidth}

\textbf{Mi compromiso para la próxima sesión es:}\\[1mm]
\linead{\linewidth}\\[1mm]
\linead{\linewidth}

\begin{softbox}{milcMostaza}{milcGris}{Cita de despedida · Marco Aurelio}
\textit{``El obstáculo en el camino se vuelve el camino. Lo que estorba al camino se vuelve camino.''}\\[1mm]
Cada error de hoy, bien observado, se convierte en pieza del camino que pisarás mañana.
\end{softbox}

\small
\begin{tabularx}{\textwidth}{>{\bfseries}p{3.6cm}Y}
\rowcolor{milcGradeDark!12}Nombre completo: & \linead{10cm}\\
Fecha: & \linead{4cm} \quad Curso/grupo: \linead{4cm}\\
Mi firma: & \linead{9cm}\\
\end{tabularx}
\normalsize

\begin{infoband}{milcGradeDark}
\textbf{Para la próxima sesión.} Dos cosas. \textbf{(1) Valida tu necesidad con su dueño:} habla con la persona a quien le serviría tu aparato (tu abuela, alguien de tu casa o tu vereda) y pregúntale si \textbf{de verdad} le ayudaría y cómo lo usaría. Anota lo que te diga ---quizás cambie tu idea, y eso es bueno: es escuchar de verdad. \textbf{(2) Observa un sistema que ya percibe y actúa:} busca en tu día un aparato que reaccione al entorno (una puerta automática, una luz con sensor, un carro que pita al retroceder) e identifica sus tres partes: \emph{¿qué percibe (sensor)? ¿qué regla decide? ¿qué hace (actuador)?}. \textbf{Trae:} tu necesidad elegida afinada (ya validada con su dueño) y el análisis de ese sistema. \emph{En el módulo 2 vas a bajar tu necesidad a pseudocódigo y diagrama de flujo ---la lógica antes que los bloques---, y en el 3 la construirás en MakeCode. El sapo ya te enseñó lo más difícil: escuchar antes de actuar.}
\end{infoband}

\end{document}
